\section{D+1午间更新：明天D+2到底怎么做}

\begin{sourcequalitybox}[先看今天已经发生的事]
本节是\textbf{9月1日12:11的午间版}，不是收盘结论。芒果超媒报20.38元，涨20.02\%；掌阅科技报24.85元、博纳影业报6.17元、中广天择报21.01元、欢瑞世纪报4.73元，均涨停。上海电影报19.03元，涨3.82\%。同时，大盘午间上涨家数多于下跌家数，但创业板、科创板下跌；市场状态仍应按防御性轮动理解，不能把今天的涨停直接当作全面进攻市。
\end{sourcequalitybox}

当天新增的公开报道明确写到《后西游记》由芒果超媒旗下芒果TV出品、芒果TV AIGC创新内容中心制作。这让芒果超媒从“只看平台页面的温度计”升级为\textbf{本项目的创业板核心观察股}；它仍是创业板，且项目收入、利润和后续数据尚未披露。主板方面，掌阅、博纳、中广天择和欢瑞的强势只能证明AI内容、影视和短剧方向在扩散，不能证明它们与《后西游记》有项目关系。\sourcenote{S12；S13。}

\begin{exhibitbox}[今天午间的梯队：强度有了，主板根基还没有]
\begin{tabularx}{\textwidth}{L{2.5cm}L{1.6cm}L{1.8cm}X}
\toprule
股票 & D+1午间表现 & 午间成交额 & 在这条题材里的正确定位 \\
\midrule
芒果超媒（300413） & +20.02\% & 24.36亿元 & 项目核心观察股；创业板，主板账户不能把它换成别的股票。 \\
掌阅科技（603533） & 涨停 & 11.63亿元 & 主板AI内容扩散股；没有项目直接关系证据。 \\
博纳影业（001330） & 涨停 & 13.80亿元 & 主板影视扩散股；没有项目直接关系证据。 \\
中广天择（603721） & 涨停 & 0.97亿元 & 主板扩散股；成交规模偏小，不能叫容量核心。 \\
欢瑞世纪（000892） & 涨停 & 0.54亿元 & 主板扩散股；不能叫项目龙头。 \\
上海电影（601595） & +3.82\% & 5.18亿元 & 影视方向跟随观察；当下不属于本项目关系股。 \\
\bottomrule
\end{tabularx}
\end{exhibitbox}

\subsection*{今天下午只做收盘确认，不新增主板“龙头”}

下午收盘后，只填三件事：第一，芒果超媒是否仍明显强于市场；第二，掌阅和博纳谁的\textbf{全天成交额}更大、收盘是否仍封在涨停；第三，有没有公司公告、交易所互动回复或项目方名单，新出现一只主板项目关系股。前两件即使都很好，也只能把题材定义为“AI内容扩散很强”；第三件不出现，就不能把题材改叫“主板AI长剧主线”。

\subsection*{明天D+2的纸面操作预案}

\textbf{情形一：收盘前仍没有主板项目关系股。} 主板账户对“AI长剧龙头”不做任何买入假设。最强的是创业板芒果超媒时，不拿掌阅、博纳或别的传媒股去替代它；这正是避免买到后排跟风的关键。

\textbf{情形二：你把范围主动放宽为“AI内容扩散”，而不是AI长剧龙头。} D+2开盘前只保留掌阅和博纳两只作比较：谁在今天收盘时仍封涨停、全天成交额更大，谁排在观察表第一。若两只都封板，就两只都不在竞价或一字涨停时排队；能买到不等于应该买。若第一名无法正常交易，也不自动换买第二名。

\textbf{情形三：第一名D+2能正常交易。} 不抢开盘。开盘后先看15分钟：价格跌破开盘价且没有收回，删掉当天的纸面参与计划；价格守住开盘价、回到涨停价并持续封住，才记为“强势延续”。即便满足，也只能把它记为\textbf{AI内容扩散股的纸面候选}，不是AI长剧龙头。因为A股T+1，D+2建立的纸面仓位最早到D+3才能处理；所以D+2不能用“当天冲高就卖”的假设去掩盖风险。

\textbf{情形四：D+2出现新的主板项目关系股首板。} 旧的扩散股全部降为陪跑，只记录新关系股。它D+2当天只是首板，仍不叫龙头；等D+3看它能否走成二板，同时是否又有新的关系股首板。只有这样，主板梯队才真正开始形成。

\subsection*{明天收盘后的一句话结论规则}

\begin{itemize}
  \item 主板项目关系股仍为零：结论写“AI内容扩散强，AI长剧主板龙头为空”。
  \item 出现一只项目关系股首板：结论写“有主板种子，等D+3确认二板”。
  \item 出现二板关系股，且新关系股首板补位：结论写“主板梯队开始形成，进入D+3高度股与容量股比较”。
  \item 芒果超媒走弱、主板扩散股大面积掉队：结论写“D+1情绪结束，删除整张观察表”。
\end{itemize}

\qualitynote{本节价格、涨幅和成交额来自本地公开行情快照，观察时间为2026年9月1日12:11，行情源未提供原始时间戳；收盘后必须重新核对。午间快照没有完整涨跌停家数和订单簿，故没有把封单金额或全市场连板高度写成结论。}
